\chapter{Background and Related Work}
\label{chap:background}

\section{Human Motion Representation}
\label{sec:bg_representation}

A motion clip is a sequence of body poses sampled at a fixed frame rate. A pose can be written as joint positions, as joint rotations along a kinematic tree, or as a redundant combination of both. The HumanML3D benchmark~\cite{guo2022humanml3d} fixes one such combination: a 263-dimensional feature vector per frame, at 20 frames per second, for a 22-joint skeleton. It contains the root's angular and linear velocity and height, the root-relative positions of the other joints, their rotations in the continuous 6D parameterization~\cite{zhou2019continuity}, the joint velocities and four binary foot-contact labels. Joint positions are recovered from this vector by integrating the root velocities and adding the root-relative positions. \chapref{chap:protocol} gives the exact layout.

Whole-body models such as SMPL-X~\cite{pavlakos2019smplx} add articulated hands and a face to the body. We use SMPL-X only to draw a body surface from the 22 recovered joints.

\section{Discrete Motion Tokenization}
\label{sec:bg_tokenization}

A motion tokenizer is an autoencoder with a discrete bottleneck. An encoder maps $T$ frames to $T/d$ latent vectors, a quantizer replaces each latent vector by one or more integer codes, and a decoder reconstructs the frames from the codes. A generator can then model motion as a sequence of integers.

\paragraph{Vector quantization.} VQ-VAE~\cite{oord2017vqvae} quantizes a latent vector $z$ to its nearest neighbour in a learned codebook $\{e_k\}_{k=1}^{V}$. The codebook is trained alongside the autoencoder and is fragile: codes that are never selected receive no gradient and stay unused. T2M-GPT~\cite{zhang2023t2mgpt} reports that a plain VQ tokenizer on HumanML3D reaches reconstruction FID $0.492$, and that exponential-moving-average codebook updates together with resetting dead codes bring this to $0.070$. A fair VQ baseline must therefore include both.

\paragraph{Residual vector quantization.} A residual quantizer applies $Q$ codebooks in sequence, each quantizing what the previous ones left over. MoMask~\cite{guo2024momask} uses this with quantizer dropout and reports reconstruction FID $0.019$.

\paragraph{Finite scalar quantization.} FSQ~\cite{mentzer2023fsq} removes the learned codebook. Each latent channel is bounded and rounded to one of a small number of levels, so the set of codes is a fixed lattice. There is nothing to collapse and nothing to reset. FSQ has been shown to outperform VQ on motion by ScaMo~\cite{lu2024scamo}, and in a multi-scale form by Wang et al.~\cite{wang2025stmsq}, who report reconstruction FID $0.037$ and generation FID $0.063$. \chapref{chap:tokenizer} positions our grouped-residual FSQ against these.

\section{Text-to-Motion Generation}
\label{sec:bg_generation}

Token-based generators fall into two groups.

\paragraph{Causal autoregressive models.} T2M-GPT~\cite{zhang2023t2mgpt} trains a decoder-only transformer to predict the next motion token given the text and the previous tokens, and reports FID $0.116$. AttT2M~\cite{zhong2023attt2m} adds body-part attention (FID $0.112$). Mogo~\cite{fu2024mogo} extends the idea to residual tokens with a hierarchical causal transformer (FID $0.079$). MotionGPT~\cite{jiang2023motiongpt} treats motion tokens as a language for a pretrained language model. These models can emit motion incrementally.

\paragraph{Masked and bidirectional models.} MoMask~\cite{guo2024momask} (FID $0.045$), MMM~\cite{pinyoanuntapong2023mmm} and BAMM~\cite{pinyoanuntapong2024bamm} predict masked tokens using context on both sides. They are the most accurate published models on HumanML3D and cannot stream, because a token is not final until the tokens after it have been considered.

\paragraph{Large motion models.} A separate line scales data and parameters: the Large Motion Model~\cite{zhang2024lmm}, the million-clip corpus and models of Wang et al.~\cite{wang2024beingm0}, and MotionMillion~\cite{fan2025motionmillion}. Scale is an axis this thesis does not explore; we cite these as motivation for the pretraining ablation of \chapref{chap:generator}.

\section{State-Space Sequence Models}
\label{sec:bg_ssm}

A linear state-space model maps an input signal $u(t)$ to an output $y(t)$ through a hidden state $x(t)$:
\begin{equation}
\dot{x}(t) = A\,x(t) + B\,u(t), \qquad y(t) = C\,x(t) + D\,u(t).
\label{eq:ssm_continuous}
\end{equation}
With a step size $\Delta$ and a zero-order hold on the input, the discrete recurrence is
\begin{equation}
s_t = \bar{A}\,s_{t-1} + \bar{B}\,u_t, \qquad \bar{A} = e^{\Delta A}, \qquad \bar{B} = (\Delta A)^{-1}\bigl(e^{\Delta A} - I\bigr)\,\Delta B .
\label{eq:ssm_discrete}
\end{equation}
S4~\cite{gu2022s4} uses such a system as a sequence layer. Mamba~\cite{gu2023mamba} makes it \emph{selective}: $\Delta$, $B$ and $C$ become functions of the current input, so the layer can decide, per token, how much of its state to keep. The layer is called S6.

Two properties matter here. During training, the recurrence can be evaluated for all positions in parallel, at a cost linear in the sequence length. During inference, one step needs only the previous state $s_{t-1}$, whose size does not depend on $t$.

Mamba has been applied to motion before, inside diffusion and masked pipelines. Examples are Motion Mamba~\cite{zhang2024motionmamba}, KMM~\cite{zhang2024kmm} and T2M Mamba~\cite{zhan2026t2mmamba}. All of these process the sequence in both directions, so none exercises the constant-state inference path.

\section{The Streaming Requirement and the Unoccupied Design Cell}
\label{sec:bg_streaming}

\begin{df}[Streaming generator]
A generator is \emph{streaming} if (i) the token emitted at step $t$ depends only on the condition and on tokens $1,\dots,t-1$, and (ii) the memory it carries between steps is bounded by a constant independent of $t$.
\end{df}

Condition (i) excludes full-sequence diffusion and bidirectional masking as the runtime model. Condition (ii) separates the two causal options. A transformer with a key--value cache stores, for each of its layers, the keys and values of all $t-1$ previous positions; its state is $O(t)$ and its per-step attention cost is $O(t)$. A state-space model stores one state per layer; both are $O(1)$.

\tabref{tab:design_cells} places prior work along three properties. Streaming systems with a continuous latent exist~\cite{xiao2025motionstreamer,li2026llamo}, and so do discrete causal transformers~\cite{zhang2023t2mgpt,fu2024mogo}. The cell that combines discrete tokens, causal next-token prediction and a fixed-size recurrent state is, to our knowledge, empty. \chapref{chap:generator} fills it.

\begin{table}[H]
\centering
\caption{Design cells for text-to-motion generators. Only the last row satisfies both streaming conditions with discrete tokens.}
\label{tab:design_cells}
\begin{tabular}{p{0.52\textwidth}ccc}
\toprule
Family & Discrete & Causal & Fixed state \\
\midrule
Masked transformers~\cite{guo2024momask,pinyoanuntapong2023mmm,pinyoanuntapong2024bamm} & yes & no & no \\
Mamba in diffusion or masked pipelines~\cite{zhang2024motionmamba,zhang2024kmm,zhan2026t2mmamba} & varies & no & no \\
Causal transformers~\cite{zhang2023t2mgpt,fu2024mogo,zhong2023attt2m} & yes & yes & no \\
Continuous-latent streaming~\cite{xiao2025motionstreamer,li2026llamo} & no & yes & no \\
Token-autoregressive state-space model (this thesis) & yes & yes & yes \\
\bottomrule
\end{tabular}
\end{table}
